% Hologram APEX — Van Gogh Pattern Recognition
% LaTeX article compiling mathematics, algorithms, and validation
\documentclass[11pt]{article}
\usepackage[a4paper,margin=1in]{geometry}
\usepackage{amsmath,amssymb,amsthm,mathtools}
\usepackage{bm}
\usepackage{graphicx}
\usepackage{booktabs}
\usepackage{algorithm}
\usepackage{algpseudocode}
\usepackage{hyperref}
\usepackage[nameinlink]{cleveref}
\usepackage{microtype}
\usepackage{enumitem}
\setlist{nosep}

\title{Phase-Aware Holographic Pattern Recognition for Van~Gogh:\\ A Low-Rank \texorpdfstring{$\Pi$}{Pi}-Kernel with Elastic Channels and Structured Sampling}
\author{Hologram APEX \;\; (prototype and analysis)}
\date{\today}

\newtheorem{theorem}{Theorem}
\newtheorem{proposition}{Proposition}
\newtheorem{lemma}{Lemma}
\newtheorem{definition}{Definition}

\newcommand{\R}{\mathbb{R}}
\newcommand{\C}{\mathbb{C}}
\newcommand{\E}{\mathbb{E}}
\newcommand{\norm}[1]{\left\lVert #1 \right\rVert}
\newcommand{\ip}[2]{\left\langle #1,\,#2 \right\rangle}
\newcommand{\bbm}[1]{\bm{#1}}

\begin{document}
\maketitle

\begin{abstract}
We formalize a phase-aware pattern recognition framework for Van~Gogh paintings grounded in a low-rank $\Pi$-kernel, elastic-channel regularization, and structured sampling. The method factors image evidence into complex multi-scale orientation harmonics, color opponency, and an impasto proxy, then performs convex coefficient refinement with FISTA. We define proximal invariants and a normalized commutator budget, prove convergence under Lipschitz conditions, and instantiate a nonnegative class-cone test for forgery screening. A complete research pipeline with ablations, interpretability, and sampling-variance analysis is specified together with a reference implementation.
\end{abstract}

\section{Introduction}
Van~Gogh's brushwork exhibits strong local orientation structure, distinctive color palettes, and tactile impasto. We target these properties with a mathematically controlled architecture: a $\Pi$-atom analysis that preserves complex phase, adapts orientation groups, and refines coefficients via proximal optimization. We expose channel weights through an elastic net rather than hard $\ell_1$ projection for capacity and stability.

\paragraph{Contributions.} (i) A low-rank $\Pi$-kernel with phase-aware channel statistics $(\mu,\kappa,\xi)$. (ii) An adaptive orientation selection rule using circular frequency. (iii) A convex objective with TV smoothness, solved by complex FISTA with guarantees. (iv) A normalized commutator budget to monitor operator coupling. (v) A deterministic six-orbit C768 sampler and a variance-ratio criterion to justify structured sampling. (vi) A complete, reproducible pipeline and CLI.

\section{Mathematical Framework}
\subsection{Ambient space and atom family}
Let $\Omega\subset\R^2$ be the image domain and $E=L^2(\Omega;\R^3)$ in Lab/opponent color. A $\Pi$-atom index is $\pi=(a,\rho,\Omega_j,k,\beta)$ with scale $j$, orientation tile $k$, and color bin $\beta$. The original projector is
\begin{equation}
R_\pi=R^{(a)}\otimes R^{(\rho)}\otimes R^{(\Omega)}_j\otimes R^{(k)}\otimes R^{(\beta)}.
\end{equation}
To control complexity we use a rank-$r$ expansion
\begin{equation}
R_\pi\approx\sum_{m=1}^{r}\lambda_{\pi,m}\; R^{(a)}_{m}\otimes R^{(\rho)}_{m}\otimes R^{(\Omega)}_{m}\otimes R^{(k)}_{m}\otimes R^{(\beta)}_{m},\quad \sum_{m\le r}\lambda_{\pi,m}^2\big/\sum_m\lambda_{\pi,m}^2\ge \tau,\;\tau\in[0.9,0.95].
\end{equation}

\subsection{Complex responses and orientation harmonics}
Given a luminance channel $I\in L^2(\Omega)$, we compute complex filter-bank responses $Z_{j,k}= (\phi_{j,k}\ast I)(x)\in\C$, using a dual-tree complex wavelet transform (DTCWT) when available or a Morlet bank otherwise. Let $K$ be the number of orientations per scale.

\paragraph{Adaptive orientation group.} Define the dominant circular frequency
\begin{equation}
 m^*\in\{6,8,12\},\quad m^*:=\arg\max_m\Bigl|\frac{1}{K}\sum_{k=0}^{K-1} e^{\,i m\,\theta_k}\Bigr|,\ \ \theta_k=\arg\bigl(\E_{x,j}[Z_{j,k}(x)]\bigr).
\end{equation}
The orientation harmonic at $m^*$ is the DFT coefficient
\begin{equation}
H_j(x)=\sum_{k=0}^{K-1}Z_{j,k}(x)\,\exp\!\left(-\,\tfrac{2\pi i}{K}\,m^*k\right),\qquad H_j: \Omega\to\C.
\end{equation}

\paragraph{Phase-aware channel statistics.} For each scale $j$ define
\begin{equation}
\mu_j=\E_x\bigl[|H_j(x)|\bigr],\quad \kappa_j=\Bigl|\E_x\bigl[e^{i\,\arg H_j(x)}\bigr]\Bigr|,\quad \xi_j=\frac{\E_x\bigl[|H_j(x)|\,e^{i\,\arg H_j(x)}\bigr]}{\E_x\bigl[|H_j(x)|\bigr]}\in\C.
\end{equation}
We concatenate $\mu,\kappa,\Re\xi,\Im\xi$ and add an impasto proxy
\begin{equation}
 r=\E_x\Bigl[\sqrt{\bigl(\partial_x G_\sigma\!\ast I\bigr)^2+\bigl(\partial_y G_\sigma\!\ast I\bigr)^2}\Bigr],\ \ \sigma\in[2,4],
\end{equation}
to obtain a compact signature.

\subsection{Coefficient refinement and invariants}
We refine harmonics $\{H_j\}_j$ as variables $X\in\C^{J\times\Omega}$ by solving
\begin{equation}\label{eq:obj}
\min_{X}\ \frac{1}{2}\norm{X-Y}_2^2 + \lambda_1\sum_{j,x} w_{j,x}|X_{j}(x)| + \frac{\lambda_2}{2}\norm{X}_2^2 + \frac{\lambda_{\mathrm{TV}}}{2}\sum_{j}\!\sum_{\langle x,x'\rangle}\!\bigl|X_j(x)-X_j(x')\bigr|^2.
\end{equation}
Here $Y$ stacks the raw harmonics. We use FISTA with complex soft-thresholding. The smooth part has Lipschitz constant $L\le 1+\lambda_2+8\lambda_{\mathrm{TV}}$ by a lattice Laplacian bound.

\begin{theorem}[Convergence]
With step $\eta\in(0,2/L)$, the forward–backward map $U=\mathrm{prox}_{\eta \lambda_1\|\cdot\|_1}\!\circ(I-\eta\nabla g)$ is nonexpansive. The Krasnosel'skii iteration $X^{t+1}=(1-\alpha)X^t+\alpha U(X^t)$ converges to the unique minimizer of \eqref{eq:obj}. FISTA achieves $O(1/t^2)$ objective decay.
\end{theorem}
\begin{proof}[Sketch]
$g$ is convex with $L$-Lipschitz gradient, and the $\ell_1$ plus quadratic TV are proper, convex, and lower semicontinuous. Standard forward–backward splitting guarantees apply; acceleration follows from Nesterov's estimate sequence.
\end{proof}

\paragraph{Proximal invariants.} The energy $E_t$ is nonincreasing; the support size $S_t=\#\{(j,x):X^t_j(x)\neq 0\}$ stabilizes after finite steps; and the normalized commutator budget
\begin{equation}
 C_t = \sum_{p<q} \frac{\norm{[\mathcal{B}_p,\mathcal{B}_q]}_F}{\norm{\mathcal{B}_p}_F\,\norm{\mathcal{B}_q}_F}
\end{equation}
is controlled by blockwise Lipschitz constants when group sparsity is active.

\subsection{Elastic channels and runtime}
Let channel groups $\{\mathcal{G}_p\}$ aggregate atoms routed to prime channels. We learn weights $w$ with an elastic-net group penalty
\begin{equation}
\min_w\ f(w)+\lambda_1\sum_p b_p\,\norm{w_{\mathcal{G}_p}}_2+\tfrac{\lambda_2}{2}\norm{w}_2^2,\qquad b_p\propto 1/\sqrt{\widehat{\mathrm{Var}}(\nabla_p f)}.
\end{equation}
A trust update $\mathcal{K}^{t+1}=(1-\gamma)\mathcal{K}^t+\gamma\sum_p w_p\mathcal{B}_p$, $\gamma\in(0,1)$, keeps contractions bounded.

\subsection{Class cone for forgery screening}
Given signatures $x\in\R^d$, and exemplars $X=[x_1,\dots,x_n]\in\R^{d\times n}$ for a class, the cone is $\mathcal{C}=\{X\alpha: \alpha\ge 0\}$. The forgery score is
\begin{equation}
\mathrm{dist}(x,\mathcal{C})=\min_{\alpha\ge 0}\norm{x-X\alpha}_2,\quad \text{solved via NNLS}.
\end{equation}

\section{Sampling and variance}
\paragraph{Stratified random baseline.} Patches are sampled on a coarse grid with quadrant stratification.

\paragraph{Deterministic C768.} We implement six interleaved low-discrepancy orbits that cover at most 768 grid cells per image. Let $\{(u_i,v_i)\}$ be Halton draws perturbed per orbit; cells are visited once producing near-uniform coverage.

\begin{definition}[Variance ratio]
For an observable $g$ over patches and an estimator $\hat\mu=\frac{1}{M}\sum_{s\in \mathcal{S}} g(s)$, define $\rho=\mathrm{Var}_{\text{Rand}}(\hat\mu) / \mathrm{Var}_{\text{C768}}(\hat\mu)$. We accept C768 if $\rho>1.2$ on repeated trials.
\end{definition}

\section{Algorithms}
\subsection{Feature extraction}
\begin{algorithm}[H]
\caption{Phase-aware feature extraction}
\begin{algorithmic}[1]
\Require image $I$, patch size $P$, scales $J$, orientations $K$
\State sample patches $\{\Omega_t\}$ (stratified or C768)
\For{each patch}
  \State compute complex responses $Z_{j,k}$
  \State select $m^*$ by circular statistic; $H_j\leftarrow\sum_k Z_{j,k}e^{-2\pi i m^*k/K}$
  \State optional FISTA on $H$ with $(\lambda_1,\lambda_2,\lambda_{\mathrm{TV}})$
  \State accumulate $\mu,\kappa,\xi$ and impasto $r$
\EndFor
\State return image-level signature by averaging patch features
\end{algorithmic}
\end{algorithm}

\subsection{FISTA refinement}
\begin{algorithm}[H]
\caption{Complex FISTA with TV}
\begin{algorithmic}[1]
\State initialize $X^0=Y$, $Z^0=X^0$, $t_0=1$, step $\eta=1/L$
\For{$t=0,1,\dots,T-1$}
  \State $G\leftarrow (Z^t-Y)+\lambda_2 Z^t+\lambda_{\mathrm{TV}}\,\Delta Z^t$
  \State $X^{t+1}\leftarrow \mathrm{soft}_{\eta\lambda_1}(Z^t-\eta G)$ \hfill (complex soft)
  \State $t_{t+1}\leftarrow \tfrac{1+\sqrt{1+4t_t^2}}{2}$, \; $Z^{t+1}\leftarrow X^{t+1}+\tfrac{t_t-1}{t_{t+1}}(X^{t+1}-X^t)$
\EndFor
\end{algorithmic}
\end{algorithm}

\section{Implementation}
A reference Python implementation provides: complex Morlet/DTCWT responses; adaptive $m^*$; phase-aware statistics; FISTA with TV; elastic-net logistic classifier; NNLS cone; deterministic C768 sampler; and utilities for ablations, feature importance, variance ratios, convergence traces, and phase coherence.

\paragraph{CLI.} \texttt{holo\_vangogh.py} exposes
\begin{itemize}
  \item \texttt{features} \; $\to$ feature extraction with \texttt{--sampling \{stratified, grid, c768\}} and FISTA controls.
  \item \texttt{train}, \texttt{eval} \; $\to$ elastic-net logistic evaluation with accuracy, macro-AUC, and ECE.
  \item \texttt{cone} \; $\to$ NNLS cone distance per sample.
\end{itemize}
Companion scripts: \texttt{run\_validation.py}, \texttt{feature\_importance.py}, \texttt{variance\_ratio.py}, \texttt{convergence\_analysis.py}, \texttt{phase\_coherence.py}.

\section{Validation Matrix}
We evaluate with 5$\times$2 cross-validation and the following ablations:
\begin{enumerate}[label=\textbf{A\arabic*}]
  \item Phase off vs. on: effect of $(\kappa,\xi)$.
  \item Proximal refinement: FISTA iterations and $(\lambda_1,\lambda_2,\lambda_{\mathrm{TV}})$.
  \item Regularization: elastic net vs. $\ell_1$.
  \item Impasto on/off.
  \item Sampling: stratified vs. C768 using variance ratio $\rho$.
\end{enumerate}
Metrics: artist-ID accuracy, macro-F1, macro-AUC, ECE, runtime per megapixel, peak RAM. Forgery screening reports cone distance and NNLS sparsity.

\section{Diagnostics}
\paragraph{Convergence analysis.} We log energy $E_t$ and provide traces to confirm smooth decay and to tune iteration counts for artistic textures.

\paragraph{Phase coherence.} Cross-scale coherence
$C=\E_x\bigl|\tfrac{1}{J}\sum_{j=1}^J e^{i\,\arg H_j(x)}\bigr|$ quantifies phase alignment; we hypothesize higher $C$ for structured Van~Gogh strokes.

\section{Limitations and Scope}
The Morlet proxy differs from DTCWT polar separability; Lab conversion and k-means bins are approximate; C768 is a deterministic heuristic. The framework is artisanal-style oriented and is not a replacement for general-purpose CNNs. We position it as a controlled, interpretable alternative.

\section{Conclusion}
We presented a compact, phase-aware $\Pi$-kernel for painterly pattern recognition with theoretical guarantees and a practical pipeline. The method is designed for rigorous study of style, brushwork, and palette in Van~Gogh, with empirical validation hooks for publication-grade evidence.

\appendix
\section{Parameter Heuristics}
Step $\eta=1.8/L$ with $L$ by power iteration; $\lambda_1=\kappa\,\mathrm{MAD}(\nabla f)$ with $\kappa\in[2,3]$; $\lambda_{\mathrm{TV}}=0.1\lambda_1$; elastic-net $(\lambda_1,\lambda_2)$ via grid on validation AUC.

\section{Pseudocode for NNLS Cone}
\begin{algorithm}[H]
\caption{NNLS cone distance}
\begin{algorithmic}[1]
\Require signature $x\in\R^d$, exemplars $X\in\R^{d\times n}$
\State solve $\min_{\alpha\ge 0}\norm{x-X\alpha}_2$ by active-set NNLS
\State return $\norm{x-X\alpha^*}_2$ and sparsity of $\alpha^*$
\end{algorithmic}
\end{algorithm}

\section{Reproducibility Commands}
\begin{verbatim}
python holo_vangogh.py features --input DATA --output feats.npz --sampling stratified
python holo_vangogh.py train --features feats.npz --model model.pkl
python holo_vangogh.py eval --features feats.npz --model model.pkl
python holo_vangogh.py cone --features feats.npz --output cone.json
python run_validation.py --input DATA --workdir results/
python feature_importance.py --model results/model_phase_fista.pkl --topk 15
python variance_ratio.py --holo holo_vangogh.py --input DATA/van_gogh/sample.jpg --mode c768
python convergence_analysis.py --holo holo_vangogh.py --input DATA/van_gogh/sample.jpg --output trace.json
python phase_coherence.py --holo holo_vangogh.py --input DATA/van_gogh/sample.jpg --output coherence.json
\end{verbatim}

\end{document}
